%% ma: L12-B02-0013
%% lop: 12
%% chuong: 1
%% bai: 2
%% dang: 12.2.9
%% loai: TN4
%% muc: TH
%% dap_an: "A"
%% nguon: "HK1-DeOn_1_1.pdf (M10, Đề ôn tập học kì 1 Toán 12, Đề số 2), Phần I Câu 6"
%% kien_thuc: "Bài 2 (giá trị lớn nhất, nhỏ nhất)"
%% trang_thai: cho-duyet
%% ly_do: "Dạng mới đề xuất 12.2.9"
%% ghi_chu: "Hình gốc chỉ cho các mốc (f(0) = 8/3, f(5) = 37/12, cực đại 8 tại x = 4, cực đại địa phương tại x = 1, cực tiểu địa phương giữa 1 và 4 nằm trên mức 37/12). Bản kho dựng đường cong bằng đa thức nội suy qua các mốc này (không theo tỉ lệ ở các điểm không ghi số)."
\begin{ex}
Cho hàm số $y=f(x)$ liên tục trên $\mathbb R$ và có đồ thị như hình bên.
\begin{center}\begin{tikzpicture}[x=0.9cm,y=0.75cm]
\draw[truc] (-0.6,0)--(6,0) node[below] {$x$};
\draw[truc] (0,-0.8)--(0,9.2) node[left] {$y$};
\node[below left] at (0,0) {$O$};
\draw[phu] (0,8)--(4,8)--(4,0) (1,0)--(1,5.5) (0,3.0833)--(5,3.0833)--(5,0);
\node[left] at (0,8) {$8$}; \node[left] at (0,3.25) {$\tfrac{37}{12}$}; \node[left] at (0,2.5) {$\tfrac{8}{3}$};
\node[below] at (1,0) {$1$}; \node[below] at (4,0) {$4$}; \node[below] at (5,0) {$5$};
\draw[thick] plot[smooth] coordinates {(-0.280,0.211) (-0.180,1.218) (-0.081,2.072) (0.019,2.795) (0.119,3.405) (0.218,3.916) (0.318,4.339) (0.417,4.685) (0.517,4.961) (0.617,5.173) (0.716,5.328) (0.816,5.431) (0.916,5.486) (1.015,5.500) (1.115,5.476) (1.215,5.422) (1.314,5.342) (1.414,5.242) (1.514,5.130) (1.613,5.010) (1.713,4.890) (1.813,4.776) (1.912,4.674) (2.012,4.591) (2.111,4.532) (2.211,4.503) (2.311,4.507) (2.410,4.548) (2.510,4.628) (2.610,4.749) (2.709,4.910) (2.809,5.109) (2.909,5.345) (3.008,5.611) (3.108,5.903) (3.207,6.212) (3.307,6.530) (3.407,6.845) (3.506,7.147) (3.606,7.422) (3.706,7.659) (3.805,7.843) (3.905,7.961) (4.005,8.000) (4.104,7.949) (4.204,7.797) (4.304,7.539) (4.403,7.169) (4.503,6.689) (4.603,6.103) (4.702,5.424) (4.802,4.673) (4.901,3.876) (5.001,3.075) (5.101,2.319) (5.200,1.675) (5.300,1.221)};
\end{tikzpicture}\end{center}
Gọi $m$, $M$ lần lượt là giá trị nhỏ nhất và giá trị lớn nhất của hàm số $f(x)$ trên đoạn $[0;5]$. Khẳng định nào sau đây là đúng?
\choice{\True $3m+M=16$}{$8m+3M=16$}{$12m-3M=13$}{$12m-4M=5$}
\loigiai{Từ đồ thị: $M=f(4)=8$; giá trị nhỏ nhất đạt tại đầu mút $x=0$: $m=f(0)=\dfrac83<f(5)=\dfrac{37}{12}$. Vậy $3m+M=8+8=16$.}
\end{ex}
